
% ISW SECTION FOR PAPER I

\subsection{Integrated Sachs-Wolfe Signature}
\label{sec:isw}

The late-time evolution of the dark energy equation of state in our model 
produces a distinctive signature in the Integrated Sachs-Wolfe (ISW) effect. 
The ISW contribution to CMB anisotropies is:
%
\begin{equation}
\left(\frac{\Delta T}{T}\right)_{\rm ISW} = 2\int_0^{z_{\rm LSS}} 
\dot{\Phi} \, \frac{dt}{dz} \, dz
\end{equation}
%
where the time derivative of the gravitational potential $\dot{\Phi}$ 
is sensitive to both the expansion history and the growth of structure.

In the Evaporating Universe, the scalar field transitions from a tracking behavior 
($w \simeq -0.05$) to a phantom-like phase ($w \simeq -1.22$) at late times. 
Our exact numerical solution of the growth equation predicts an **enhanced** ISW signal, 
with a peak deviation of approximately **6.8\%** relative to $\Lambda$CDM at $z \simeq 1.25$.

While current ISW detections ($\sim 3\sigma$ with Planck \cite{Planck2018}) 
are insufficient to distinguish this signal, upcoming surveys offer promising 
prospects. Cross-correlations between Euclid galaxy samples and Planck CMB maps, 
or LSST $\times$ CMB-S4, could achieve the sensitivity needed to test this 
prediction. This provides a \textit{falsifiable} test of the model's 
late-time dynamics.
